%% Đề 10 -- de-giua-ky-1-toan-9-nam-2022-2023-phong-gddt-huu-lung-lang-son.pdf
\deheader{ĐỀ SỐ 10}%
  {Đề kiểm tra Giữa Học kỳ — Năm học 2022--2023}%
  {Phòng GD\&ĐT | Chỉ phần Đại số}
\addcontentsline{toc}{section}{Đề 10 — Đề kiểm tra Giữa HK — Năm học 2022--2023}

\begin{de_cau}{Câu 1}{2,0 điểm}
\begin{enumerate}[label=\alph*)]
  \item Tính giá trị biểu thức: $A = \sqrt{(1-\sqrt{5})^2} + \sqrt{5}$.
  \item Rút gọn biểu thức: $B = \dfrac{x-4}{\sqrt{x}-2}$ với $x \ge 0, x \neq 4$.
\end{enumerate}
\end{de_cau}

\begin{loi_giai_cau}
\textbf{a)} Ta có:
\[A = \sqrt{(1-\sqrt{5})^2} + \sqrt{5} = |1-\sqrt{5}| + \sqrt{5}.\]
Vì $1 < \sqrt{5}$ nên $1 - \sqrt{5} < 0$, suy ra $|1-\sqrt{5}| = \sqrt{5} - 1$.

Do đó:
\[A = \sqrt{5} - 1 + \sqrt{5} = \boxed{2\sqrt{5} - 1}.\]

\textbf{b)} Với $x \ge 0, x \neq 4$, ta có:
\[B = \dfrac{x-4}{\sqrt{x}-2} = \dfrac{(\sqrt{x}-2)(\sqrt{x}+2)}{\sqrt{x}-2} = \boxed{\sqrt{x}+2}.\]
\end{loi_giai_cau}

\begin{de_cau}{Câu 2}{2,5 điểm}
\begin{enumerate}[label=\alph*)]
  \item Giải phương trình: $2x^2 - 3x - 2 = 0$.
  \item Giải hệ phương trình:
  \[\begin{cases} 3x + y = 7 \\ x - 2y = -7 \end{cases}.\]
  \item Cho hàm số $y = ax^2$. Tìm hệ số $a$ biết đồ thị hàm số đi qua điểm $M(-3; 9)$.
\end{enumerate}
\end{de_cau}

\begin{loi_giai_cau}
\textbf{a)} Phương trình $2x^2 - 3x - 2 = 0$ có $a=2, b=-3, c=-2$.

Biệt thức:
\[\Delta = b^2 - 4ac = (-3)^2 - 4\cdot 2\cdot (-2) = 9 + 16 = 25 > 0.\]
Vì $\Delta > 0$ nên phương trình có hai nghiệm phân biệt:
\[x_1 = \dfrac{-b+\sqrt{\Delta}}{2a} = \dfrac{3+5}{4} = 2, \quad x_2 = \dfrac{-b-\sqrt{\Delta}}{2a} = \dfrac{3-5}{4} = -\dfrac{1}{2}.\]
Vậy tập nghiệm của phương trình là $S = \left\{2; -\dfrac{1}{2}\right\}$.

\textbf{b)} Nhân cả hai vế của phương trình thứ nhất với $2$:
\[\begin{cases} 3x + y = 7 \\ x - 2y = -7 \end{cases} \Leftrightarrow \begin{cases} 6x + 2y = 14 \\ x - 2y = -7 \end{cases}.\]
Cộng từng vế hai phương trình ta được:
\[7x = 7 \Leftrightarrow x = 1.\]
Thay $x = 1$ vào phương trình $3x + y = 7$, ta có:
\[3\cdot 1 + y = 7 \Leftrightarrow y = 4.\]
Vậy hệ phương trình có nghiệm duy nhất $\boxed{(x; y) = (1; 4)}$.

\textbf{c)} Đồ thị hàm số $y = ax^2$ đi qua điểm $M(-3; 9)$ khi và chỉ khi:
\[9 = a\cdot (-3)^2 \Leftrightarrow 9 = 9a \Leftrightarrow \boxed{a = 1}.\]
Vậy hàm số cần tìm là $y = x^2$.
\end{loi_giai_cau}

\begin{de_cau}{Câu 3}{2,0 điểm}
Một lớp học có 40 học sinh. Nếu xếp thành từng nhóm $n$ người thì được một số nhóm và thừa 4 người. Nếu xếp thành từng nhóm $n+2$ người thì cũng được số nhóm như thế nhưng thiếu 2 người. Tìm $n$ ($n \in \mathbb{N}^*$).
\end{de_cau}

\begin{loi_giai_cau}
Gọi số nhóm học sinh được chia là $k$ ($k \in \mathbb{N}^*$).

Theo đề bài:
\begin{itemize}
  \item Khi xếp thành từng nhóm $n$ người ($n > 4$) thì thừa 4 người, nên ta có:
  \[40 = k\cdot n + 4 \Rightarrow k\cdot n = 36.\]
  \item Khi xếp thành từng nhóm $n+2$ người thì thiếu 2 người, nên ta có:
  \[40 = k(n+2) - 2 \Rightarrow k(n+2) = 42.\]
\end{itemize}
Trừ từng vế hai phương trình:
\[k(n+2) - kn = 42 - 36 \Leftrightarrow 2k = 6 \Leftrightarrow k = 3 \text{ (thỏa mãn)}.\tag{nhận}\]
Thay $k = 3$ vào $k\cdot n = 36$, ta được:
\[3n = 36 \Leftrightarrow n = 12 \text{ (thỏa mãn } n > 4).\]
Kiểm tra lại:
\begin{itemize}
  \item Chia thành 3 nhóm mỗi nhóm 12 người: $3 \times 12 + 4 = 40$ (thừa 4 người).
  \item Chia thành 3 nhóm mỗi nhóm 14 người: $3 \times 14 - 2 = 40$ (thiếu 2 người).
\end{itemize}
Vậy $\boxed{n = 12}$.
\end{loi_giai_cau}

\begin{de_cau}{Câu 4}{0,5 điểm}
Tìm giá trị nhỏ nhất của biểu thức:
\[P = x^2 - 4x + y^2 + 2y + 6.\]
\end{de_cau}

\begin{loi_giai_cau}
Biến đổi biểu thức $P$:
\begin{align*}
P &= (x^2 - 4x + 4) + (y^2 + 2y + 1) + 1 \\
  &= (x - 2)^2 + (y + 1)^2 + 1.
\end{align*}
Vì $(x - 2)^2 \ge 0$ và $(y + 1)^2 \ge 0$ với mọi $x, y \in \mathbb{R}$, nên:
\[P \ge 0 + 0 + 1 = 1.\]
Đẳng thức xảy ra khi và chỉ khi:
\[\begin{cases} x - 2 = 0 \\ y + 1 = 0 \end{cases} \Leftrightarrow \begin{cases} x = 2 \\ y = -1 \end{cases}.\]
Vậy giá trị nhỏ nhất của $P$ là $\boxed{P_{\min} = 1}$, đạt được khi $x = 2$ và $y = -1$.
\end{loi_giai_cau}
